% =============================================================================
% Appendix: Code Migration — MATLAB to C++
% =============================================================================
\newpage
\section{Code Migration: MATLAB to C++}
\label{sec:migration}

\subsection{Motivation}
\label{subsec:migration-motivation}

The original tracking control simulation was implemented as a collection of
MATLAB scripts (\texttt{Fig2a.m}--\texttt{Fig2c.m}),
following the reference implementation in~\cite{tymoshchuk2024tracking}.
While MATLAB provides rapid prototyping through built-in matrix operations and
plotting, several factors motivated a migration to C++:

\begin{enumerate}
	\item \textbf{Execution speed.} The simulation loop iterates 45{,}000 time
	      steps with per-step nonlinear dynamics evaluation. C++ compiled code
	      with direct memory access eliminates the interpreter overhead inherent
	      to MATLAB's execution model.
	\item \textbf{Memory control.} C++ provides explicit control over stack
	      versus heap allocation, cache-aligned data layouts, and
	      zero-copy data access through pointers---critical when storing
	      $2 \times 45{,}000$ state matrices.
\end{enumerate}

\subsection{Architecture Decisions}
\label{subsec:migration-architecture}

The MATLAB implementation uses a procedural style: separate 1-D arrays for
each state variable (\texttt{x}, \texttt{y}, \texttt{r1}, \texttt{r2},
\texttt{e1}, \texttt{e2}, \texttt{u}) with a single \texttt{for} loop
driving the simulation. The C++ port uses 2d arrays for the state variables
with a single for loop as well. It adopts a struct-based architecture that
groups related data and separates concerns across compilation units.

\subsubsection{Data Layout}
\begin{table}[htbp]
	\centering
	\caption{Data structure mapping between MATLAB and C++ implementations.}
	\label{tab:data-layout}
	\begin{tabular}{lll}
		\toprule
		MATLAB Variable                        & C++ Equivalent                                       & Type                                \\
		\midrule
		\texttt{x(1:n)}, \texttt{y(1:n)}       & \texttt{stateMatrix}                                 & \texttt{Matrix<double, 2, Dynamic>} \\
		\texttt{r1(1:n)}, \texttt{r2(1:n)}     & \texttt{referenceMatrix}                             & \texttt{Matrix<double, 2, Dynamic>} \\
		\texttt{e1(1:n)}, \texttt{e2(1:n)}     & \texttt{x1ErrorMatrix}, \texttt{x2ErrorMatrix}       & \texttt{Matrix<double, 2, Dynamic>} \\
		\texttt{r2(i+1)-r2(i)}                 & \texttt{deltaReference}                              & \texttt{Matrix<double, 2, Dynamic>} \\
		\texttt{alpha}, \texttt{beta}, \ldots  & \texttt{InternalConditions}                          & \texttt{struct}                     \\
		\texttt{n}, \texttt{tmax}, \texttt{dt} & \texttt{SimSettings}                                 & \texttt{struct}                     \\
		\bottomrule
	\end{tabular}
\end{table}

The key design choice is consolidating per-variable 1-D arrays into
$2 \times N$ Eigen matrices. Row~0 stores the conversion state
$x_1$, and row~1 stores the temperature state $x_2$. This layout
preserves temporal locality: accessing both states at time step $k$ requires
reading two adjacent memory locations in the same column.

\subsubsection{Build System}

The C++ project uses CMake 3.15+ with \texttt{FetchContent} to manage two
external dependencies:

\begin{table}[htbp]
	\centering
	\caption{C++ build dependencies fetched at configure time.}
	\label{tab:dependencies}
	\begin{tabular}{llll}
		\toprule
		Library & Source                & Version & Purpose                       \\
		\midrule
		Eigen   & GitLab                & master  & Matrix algebra                \\
		cxxplot & US Naval Research Lab & v0.4.1  & Interactive Qt-based plotting \\
		\bottomrule
	\end{tabular}
\end{table}

The build system produces three executables from separate translation units:
\texttt{fig2a}, \texttt{fig2b}, \texttt{fig2c} (CSTR system). Each executable contains
its own \texttt{main()} function, with shared simulation logic factored into
\texttt{state\_space.cpp}.

\subsection{Implementation --- CSTR System (Fig.~2)}
\label{subsec:migration-implementation}

This section presents side-by-side comparisons of the MATLAB and C++
implementations for each major component of the CSTR simulation.

\subsubsection{System Dynamics}
\label{subsubsec:dynamics}

The CSTR system has two states: conversion $x_1$ and temperature $x_2$.
Since the control gain $g_1 = 0$, the conversion state evolves under
open-loop natural dynamics. The temperature state uses feedback-linearizing
sliding-mode control.

\paragraph{Conversion state $x_1$ --- open loop.}
The natural dynamics follow the discretized CSTR model:
\begin{equation}
	x_1(k{+}1) = x_1(k) + \tau\!\left[
		-\alpha\, x_1(k) + D_a\bigl(1 - x_1(k)\bigr)
		\exp\!\left(\frac{x_2(k)}{1 + x_2(k)/\gamma}\right)
		\right]
	\label{eq:x1-dynamics}
\end{equation}

\noindent
\textbf{MATLAB} (\texttt{Fig2a.m}, line~34):
\begin{lstlisting}[style=matlab]
f1(i) = x(i) + dt*(-alpha*x(i) + Da*(1-x(i))*exp(y(i)/(1+y(i)/gamma)));
\end{lstlisting}

\noindent
\textbf{C++} (\texttt{state\_space.cpp}, lines~23--28):
\begin{lstlisting}[style=cpp]
stateConditions.stateMatrix(0, i) =
    *prevX1 +
    settings.tau *
        (-internalConditions.alpha * *prevX1 +
         internalConditions.Da * (1 - *prevX1) *
             exp(*prevX2 / (1 + *prevX2 / internalConditions.gamma)));
\end{lstlisting}

The C++ version accesses parameters through the \texttt{InternalConditions}
struct rather than workspace variables, and uses cached pointer variables
(\texttt{prevX1}, \texttt{prevX2}) to avoid repeated matrix
indexing within the inner loop.

\paragraph{Temperature state $x_2$ --- closed loop.}
The feedback-linearizing control law cancels the nonlinear dynamics $f_2(x)$,
yielding the simplified closed-loop update:
\begin{equation}
	x_2(k{+}1) = x_2(k) + \tau\!\left[
	-\beta\,\operatorname{sgn}(e_2(k)) + \Delta r_2(k)
	\right]
	\label{eq:x2-closedloop}
\end{equation}
where $e_2(k) = x_2(k) - r_2(k)$ is the tracking error and
$\Delta r_2(k) = r_2(k{+}1) - r_2(k)$ is the reference change.

\noindent
\textbf{MATLAB} (\texttt{Fig2a.m}, line~67):
\begin{lstlisting}[style=matlab]
y(i+1) = y(i) + dt*(-betta*S(i) + r2(i+1) - r2(i));
\end{lstlisting}

\noindent
\textbf{C++} (\texttt{state\_space.cpp}, lines~31--34):
\begin{lstlisting}[style=cpp]
deltaR2 = &stateConditions.deltaReference(1, i - 1);
stateConditions.stateMatrix(1, i) =
    *prevX2 +
    settings.tau * (-internalConditions.beta * sgn(e2) + *deltaR2);
\end{lstlisting}

A notable difference is that the C++ implementation pre-computes the
$\Delta r_2$ matrix in a single Eigen operation before entering the simulation
loop, whereas MATLAB computes \texttt{r2(i+1)-r2(i)} inline at each step.

\subsubsection{Reference Trajectory}
\label{subsubsec:reference}

The piecewise-constant reference switches between two operating points at
$t = 15$\,s and $t = 30$\,s.

\noindent
\textbf{MATLAB} uses time-based conditionals:
\begin{lstlisting}[style=matlab]
if 0<=t(i+1) & t(i+1)<=15
    r1(i+1) = 0.4472;
elseif 15<t(i+1) & t(i+1)<30
    r1(i+1) = 0.7646;
elseif 30<=t(i+1) & t(i+1)<tmax
    r1(i+1) = 0.4472;
end
\end{lstlisting}

\noindent
\textbf{C++} uses index-based ternary expressions:
\begin{lstlisting}[style=cpp]
stateConditions.referenceMatrix(0, i) =
    i > 15000 && i < 30000 ? 0.7646 : 0.4472;
\end{lstlisting}

The index-based approach is valid because $t = i \cdot \tau$ and
$\tau = 45/45000 = 0.001$\,s, so $t = 15$ corresponds to $i = 15{,}000$.
Additionally, the C++ version pre-computes the reference difference matrix
$\Delta R$ via Eigen column operations:
\begin{lstlisting}[style=cpp]
stateConditions.deltaReference =
    stateConditions.referenceMatrix.rightCols(settings.timeSteps - 1) -
    stateConditions.referenceMatrix.leftCols(settings.timeSteps - 1);
\end{lstlisting}

\subsubsection{Signum Function}
\label{subsubsec:signum}

The signum function $\operatorname{sgn}(\cdot)$ is central to the
sliding-mode controller.

\noindent
\textbf{MATLAB} implements it with explicit branching:
\begin{lstlisting}[style=matlab]
if (y(i)-r2(i)) > 0
    S(i) = 1;
elseif (y(i)-r2(i)) < 0
    S(i) = -1;
else
    S(i) = 0;
end
\end{lstlisting}

\noindent
\textbf{C++} uses a branchless template:
\begin{lstlisting}[style=cpp]
template <typename T>
int sgn(const T &val) {
    return (T(0) < val) - (val < T(0));
}
\end{lstlisting}

The C++ implementation avoids conditional branching by exploiting boolean
arithmetic: the expression evaluates to $+1$, $-1$, or $0$ without any
\texttt{if}/\texttt{else} statements, which can improve pipeline utilization
on modern processors.

\subsubsection{Visualization}
\label{subsubsec:visualization}

MATLAB's native \texttt{subplot}/\texttt{plot} functions were replaced with
the \texttt{cxxplot} library from the US Naval Research Laboratory. The
library provides interactive Qt-based plot windows with a named-parameter API:

\begin{lstlisting}[style=cpp]
auto x1Plot = cxxplot::plot(tX1, x1Ds,
    window_title_ = "Conversion",
    show_legend_  = true,
    name_         = "x1",
    xlabel_       = "Time (s)",
    ylabel_       = "x1, r1");
x1Plot.add_graph(tX1Ref, x1RefDs, name_ = "r1");
\end{lstlisting}

To handle the 45{,}000 data points efficiently, a min/max envelope
downsampling utility reduces the data to approximately 2{,}000 points
while preserving visual fidelity of peaks and transients.

\autoref{tab:plot-mapping} maps the MATLAB plotting calls to their C++
equivalents.

\begin{table}[htbp]
	\centering
	\caption{Visualization API mapping.}
	\label{tab:plot-mapping}
	\begin{tabular}{ll}
		\toprule
		MATLAB                          & C++ (cxxplot)                                         \\
		\midrule
		\texttt{plot(t, x, 'k-')}       & \texttt{cxxplot::plot(t, x)}                          \\
		\texttt{xlabel('Time (s)')}     & \texttt{xlabel\_ = "Time (s)"}                        \\
		\texttt{ylabel('x\_1, r\_1')}   & \texttt{ylabel\_ = "x1, r1"}                          \\
		\texttt{legend('x\_1', 'r\_1')} & \texttt{show\_legend\_ = true}, \texttt{name\_ = ...} \\
		\texttt{subplot(2,1,k)}         & separate \texttt{cxxplot::plot()} call per window     \\
		\bottomrule
	\end{tabular}
\end{table}

\subsection{Performance Comparison}
\label{subsec:migration-performance}

The C++ simulation includes instrumentation via
\texttt{std::chrono::high\_resolution\_clock}. The MATLAB simulation uses
\texttt{tic}/\texttt{toc}. Both implementations were executed 20 times and
averaged. \autoref{fig:fig2-exec-time}
compares execution times across implementations.

\begin{figure}[htbp]
	\centering
	\includesvg[width=\textwidth]{fig2_exec_time_comparison}
	\caption{Execution time comparison for the CSTR tracking control
		simulation ($n = 45{,}000$ steps, $\tau = 0.001$\,s).
		Average over 20 runs.}
	\label{fig:fig2-exec-time}
\end{figure}

Note that the MATLAB averages include a first-run JIT compilation overhead,
which is representative of interactive usage. Excluding the warmup run, the
steady-state MATLAB time is approximately 0.005\,s, reducing the speedup
factor to approximately 4$\times$. The JIT overhead is a real cost in the
MATLAB execution model and is therefore included in the reported averages.

\subsubsection{Output Comparison --- CSTR System}

\autoref{fig:fig2-plot-comparison} presents side-by-side output from both
implementations to verify numerical equivalence.

\begin{figure}[htbp]
	\centering
	\begin{subfigure}[b]{0.65\textwidth}
		\centering
		\includegraphics[width=\textwidth]{matlab_fig2a_tracking}
		\caption{MATLAB (\texttt{Fig2a.m}).}
		\label{fig:fig2-matlab-tracking}
	\end{subfigure}
	\par\vspace{0.8em}
	\begin{subfigure}[b]{0.48\textwidth}
		\centering
		\includegraphics[width=\textwidth]{cxxplot_time_vs_x1_r1}
		\caption{C++ (\texttt{fig2a}, conversion $x_1$).}
		\label{fig:fig2-cpp-x1}
	\end{subfigure}
	\hfill
	\begin{subfigure}[b]{0.48\textwidth}
		\centering
		\includegraphics[width=\textwidth]{cxxplot_time_vs_x2_r2}
		\caption{C++ (\texttt{fig2a}, temperature $x_2$).}
		\label{fig:fig2-cpp-x2}
	\end{subfigure}
	\caption{Comparison of CSTR state trajectories from the MATLAB reference
		implementation and the C++ port. Both simulations use identical
		parameters ($\alpha = 1$, $\beta = 100$, $\lambda = 0.3$,
		$\gamma = 20$, $D_a = 0.072$) and initial conditions
		$x(0) = \mathbf{0}$.}
	\label{fig:fig2-plot-comparison}
\end{figure}

\subsection{Output Validation}
\label{subsec:migration-validation}

Verifying numerical equivalence between two independent implementations
requires more than visual comparison of plots. The migration includes a
CSV-based validation pipeline that compares the C++ simulation output against
MATLAB reference data at every time step.

\subsubsection{Validation Pipeline}

Both implementations export identical 7-column CSV files containing
$\bigl(t,\; x_1,\; x_2,\; r_1,\; r_2,\; e_1,\; e_2\bigr)$
at each of the $n = 45{,}000$ time steps. The MATLAB script writes its
reference data via \texttt{writematrix} at the end of the simulation loop;
the C++ executable writes its output through \texttt{std::ofstream} with
15-digit precision.

A standalone validation executable (\texttt{validate\_output}) loads both CSV
files, extracts corresponding columns into Eigen vectors, and computes the
element-wise maximum absolute error for each variable:
\begin{equation}
	\varepsilon_{\max}^{(v)} = \max_{k=0}^{n-1}
	\bigl| v_{\text{C++}}(k) - v_{\text{MATLAB}}(k) \bigr|,
	\quad v \in \{x_1, x_2, e_1, e_2\}
	\label{eq:max-abs-error}
\end{equation}

Each variable passes if $\varepsilon_{\max}^{(v)} < \epsilon$ for a
configurable tolerance $\epsilon$. The validator reports per-variable
pass/fail status and returns a nonzero exit code on any failure, enabling
integration into automated build pipelines.

\subsubsection{Validation Results}

After correcting an off-by-one initialization error in the reference
trajectory (see~\S\ref{subsubsec:reference}), the validation confirms
machine-precision agreement between the two implementations:

\begin{table}[htbp]
	\centering
	\caption{Maximum absolute error between C++ and MATLAB outputs
		for the CSTR tracking control simulation
		($n = 45{,}000$, $\tau = 0.001$\,s).}
	\label{tab:validation-results}
	\begin{tabular}{lcl}
		\toprule
		Variable                & $\varepsilon_{\max}$       & Status \\
		\midrule
		$x_1$ (conversion)     & $\sim 10^{-15}$ & PASS   \\
		$x_2$ (temperature)    & $\sim 10^{-15}$ & PASS   \\
		$e_1$ (tracking error) & $\sim 10^{-15}$ & PASS   \\
		$e_2$ (tracking error) & $\sim 10^{-15}$ & PASS   \\
		\bottomrule
	\end{tabular}
\end{table}

The errors at $\mathcal{O}(10^{-15})$ are consistent with IEEE~754
double-precision machine epsilon ($\approx 2.22 \times 10^{-16}$),
confirming that the two implementations compute identical floating-point
results up to rounding in the final digits. The default tolerance of
$\epsilon = 10^{-6}$ provides a margin of nine orders of magnitude above
the observed error floor.

\subsection{Remaining Work}
\label{subsec:migration-remaining}

The following component has not yet been migrated from MATLAB to C++:

\begin{enumerate}
	\item \textbf{Adversarial attack simulations.}
	      PGD attacks on the temperature measurement channel
	      and neural network approximation-based attacks (planned
	      implementations using IBM ART).
\end{enumerate}
